\noindent
\textbf{\large LISTA DE SÍMBOLOS Y ACRÓNIMOS}

\vspace{0.4cm}

% ============================================================
% ACRÓNIMOS Y ABREVIATURAS
% ============================================================
\noindent\textbf{Acrónimos y Abreviaturas}

\vspace{0.25cm}

\noindent
\begin{longtable}{p{3.2cm} p{10.8cm}}
\toprule
\textbf{Acrónimo} & \textbf{Significado} \\
\midrule
\endfirsthead
\multicolumn{2}{l}{\small\textit{(Continuación)}} \\[2pt]
\toprule
\textbf{Acrónimo} & \textbf{Significado} \\
\midrule
\endhead
\midrule
\multicolumn{2}{r}{\small\textit{Continúa en la página siguiente}} \\
\endfoot
\bottomrule
\endlastfoot
AIC     & \textit{Akaike Information Criterion} — Criterio de Información de Akaike \\[3pt]
ANN     & \textit{Artificial Neural Network} — Red Neuronal Artificial \\[3pt]
AR      & \textit{AutoRegressive} — Autorregresivo \\[3pt]
ARIMA   & \textit{AutoRegressive Integrated Moving Average} — Media Móvil Integrada Autorregresiva \\[3pt]
BIC     & \textit{Bayesian Information Criterion} — Criterio de Información Bayesiano \\[3pt]
BPTT    & \textit{Backpropagation Through Time} — Retropropagación a Través del Tiempo \\[3pt]
COVID-19 & \textit{Coronavirus Disease 2019} — Enfermedad por Coronavirus 2019 \\[3pt]
DL      & \textit{Deep Learning} — Aprendizaje Profundo \\[3pt]
DMH     & Dirección de Meteorología e Hidrología \\[3pt]
GPU     & \textit{Graphics Processing Unit} — Unidad de Procesamiento Gráfico \\[3pt]
GRU     & \textit{Gated Recurrent Unit} — Unidad Recurrente con Compuertas \\[3pt]
INC     & Industria Nacional del Cemento \\[3pt]
KGE     & \textit{Kling-Gupta Efficiency} — Eficiencia de Kling-Gupta \\[3pt]
LR      & \textit{Learning Rate} — Tasa de Aprendizaje \\[3pt]
LSTM    & \textit{Long Short-Term Memory} — Memoria de Corto y Largo Plazo \\[3pt]
MA      & \textit{Moving Average} — Media Móvil \\[3pt]
ML      & \textit{Machine Learning} — Aprendizaje Automático \\[3pt]
NSE     & \textit{Nash-Sutcliffe Efficiency} — Eficiencia de Nash-Sutcliffe \\[3pt]
OMM     & Organización Meteorológica Mundial \\[3pt]
RMSE    & \textit{Root Mean Squared Error} — Raíz del Error Cuadrático Medio \\[3pt]
RNN     & \textit{Recurrent Neural Network} — Red Neuronal Recurrente \\[3pt]
SARIMA  & \textit{Seasonal AutoRegressive Integrated Moving Average} \\[3pt]
SARIMAX & \textit{Seasonal AutoRegressive Integrated Moving Average with eXogenous inputs} \\[3pt]
TPE     & \textit{Tree-structured Parzen Estimator} — Estimador de Parzen con Estructura de Árbol \\[3pt]
\end{longtable}

\vspace{0.5cm}

% ============================================================
% SÍMBOLOS — LETRAS LATINAS
% ============================================================
\noindent\textbf{Símbolos — Letras Latinas}

\vspace{0.25cm}

\noindent
\begin{longtable}{p{3.5cm} p{10.5cm}}
\toprule
\textbf{Símbolo} & \textbf{Descripción} \\
\midrule
\endfirsthead
\multicolumn{2}{l}{\small\textit{(Continuación)}} \\[2pt]
\toprule
\textbf{Símbolo} & \textbf{Descripción} \\
\midrule
\endhead
\midrule
\multicolumn{2}{r}{\small\textit{Continúa en la página siguiente}} \\
\endfoot
\bottomrule
\endlastfoot
$B$                      & Operador de rezago (\textit{backshift}): $BY_t = Y_{t-1}$ \\[3pt]
$b_i,\,b_f,\,b_o$        & Vectores de sesgo de las compuertas de entrada, olvido y salida en LSTM \\[3pt]
$b_r,\,b_z,\,b_h$        & Vectores de sesgo de las compuertas de reinicio, actualización y estado candidato en GRU \\[3pt]
$C_t$                    & Estado interno (celda) de la red LSTM en el instante $t$ \\[3pt]
$\tilde{C}_t$            & Estado candidato de la celda LSTM en el instante $t$ \\[3pt]
$d$                      & Orden de diferenciación no estacional en modelos ARIMA/SARIMAX \\[3pt]
$D$                      & Orden de diferenciación estacional en modelos SARIMA/SARIMAX \\[3pt]
$F_t$                    & Compuerta de olvido (\textit{forget gate}) de la celda LSTM \\[3pt]
$H_t$                    & Estado oculto de la red recurrente en el instante $t$ \\[3pt]
$\tilde{H}_t$            & Estado oculto candidato de la celda GRU en el instante $t$ \\[3pt]
$I_t$                    & Compuerta de entrada (\textit{input gate}) de la celda LSTM \\[3pt]
$n$                      & Número total de observaciones \\[3pt]
$O_t$                    & Compuerta de salida (\textit{output gate}) de la celda LSTM \\[3pt]
$p$                      & Orden del componente autorregresivo (AR) no estacional \\[3pt]
$P$                      & Orden del componente autorregresivo estacional (SAR) \\[3pt]
$q$                      & Orden del componente de media móvil (MA) no estacional \\[3pt]
$Q$                      & Orden del componente de media móvil estacional (SMA) \\[3pt]
$R_t$                    & Compuerta de reinicio (\textit{reset gate}) de la celda GRU \\[3pt]
$R^{2}$                  & Coeficiente de determinación \\[3pt]
$s$                      & Longitud del ciclo estacional (e.g.\ $s=12$ para datos mensuales) \\[3pt]
$W_{ji}$                 & Peso sináptico de la conexión entre la neurona $i$ y la neurona $j$ \\[3pt]
$W_{xi},\,W_{xf},\,W_{xo}$ & Matrices de pesos de la entrada en las compuertas LSTM \\[3pt]
$W_{hi},\,W_{hf},\,W_{ho}$ & Matrices de pesos del estado oculto previo en las compuertas LSTM \\[3pt]
$W_{xr},\,W_{xz},\,W_{xh}$ & Matrices de pesos de entrada en las compuertas GRU \\[3pt]
$W_{hr},\,W_{hz},\,W_{hh}$ & Matrices de pesos del estado oculto en las compuertas GRU \\[3pt]
$x_n$                    & Variables de entrada de la red neuronal \\[3pt]
$X_t$                    & Vector de variables exógenas observadas en el instante $t$ \\[3pt]
$y_i$                    & Valor observado (real) de la variable objetivo \\[3pt]
$\hat{y}_i$              & Valor predicho por el modelo \\[3pt]
$Y_t$                    & Variable endógena (serie temporal objetivo) en el instante $t$ \\[3pt]
$Z_t$                    & Compuerta de actualización (\textit{update gate}) de la celda GRU \\[3pt]
\end{longtable}

\vspace{0.5cm}

% ============================================================
% SÍMBOLOS — LETRAS GRIEGAS Y OPERADORES
% ============================================================
\noindent\textbf{Símbolos — Letras Griegas y Operadores Especiales}

\vspace{0.25cm}

\noindent
\begin{longtable}{p{3.5cm} p{10.5cm}}
\toprule
\textbf{Símbolo} & \textbf{Descripción} \\
\midrule
\endfirsthead
\multicolumn{2}{l}{\small\textit{(Continuación)}} \\[2pt]
\toprule
\textbf{Símbolo} & \textbf{Descripción} \\
\midrule
\endhead
\midrule
\multicolumn{2}{r}{\small\textit{Continúa en la página siguiente}} \\
\endfoot
\bottomrule
\endlastfoot
$\alpha_i$               & Coeficientes autorregresivos del modelo AR \\[3pt]
$\beta$                  & Vector de coeficientes asociados a las variables exógenas en SARIMAX \\[3pt]
$\varepsilon_t$          & Término de error o ruido blanco en el instante $t$;\; $\varepsilon_t \sim \mathcal{N}(0,\sigma^2)$ \\[3pt]
$\theta_j$               & Coeficientes del componente de media móvil (MA) \\[3pt]
$\phi_p(B)$              & Polinomio autorregresivo no estacional \\[3pt]
$\Phi_P(B^s)$            & Polinomio autorregresivo estacional \\[3pt]
$\theta_q(B)$            & Polinomio de media móvil no estacional \\[3pt]
$\Theta_Q(B^s)$          & Polinomio de media móvil estacional \\[3pt]
$\sigma^2$               & Varianza del término de error del modelo \\[3pt]
$\sigma(\cdot)$          & Función sigmoide:\; $\sigma(x) = \frac{1}{1+e^{-x}}$ \\[3pt]
$\tanh(\cdot)$           & Función tangente hiperbólica:\; $\tanh(x) = \frac{e^x - e^{-x}}{e^x + e^{-x}}$ \\[3pt]
$\varphi$                & Función de activación genérica de una neurona artificial \\[3pt]
$\nabla^d$               & Operador de diferencia de orden $d$:\; $\nabla^d Y_t = (1-B)^d Y_t$ \\[3pt]
$\odot$                  & Producto de Hadamard (multiplicación elemento a elemento) \\[3pt]
$\mathbb{R}^{n \times h}$ & Espacio de matrices reales de dimensión $n \times h$ \\[3pt]
\end{longtable}
